% Chapter 1 -- Tokens and constants (model: TP 6.0 PG Ch 1, [bor s4.2]).
% Sources (Law 1): os/tps/tps.pas SkipTrivia, ScanIdentifier, ScanNumber,
% ScanString, NextToken, KeywordKind, IsSpace, IsIdentStart/IsIdentChar,
% FailLex; os/tps/LEXER-DUMP.md; ADR-0007 DEC-02 (case-insensitive
% identifiers, 32-bit integer).  The reserved-word table registers every
% word of the shipped lexer (tps_gate inventory, KIND reserved); the
% lexical diagrams are grammar/lexical.ebnf.  Examples replayed by
% test-pubs-tps.
\Gates{test-tps-lex-fpc, test-tps-lex-os, test-compiler, test-pubs-tps}
\chapter{Tokens and constants}
\ix{tokens}

A Turbo Initech program is a sequence of characters that the compiler
reads from left to right and breaks into \term{tokens}: special symbols,
reserved words, identifiers, numbers and character strings.  Tokens are
separated by blanks and comments.  This chapter defines each kind of token
and explains the syntax diagrams used throughout this book.

\section{Reading the syntax diagrams}
\ix{syntax diagrams}
The grammar of Turbo Initech is given by syntax diagrams.  To read a
diagram, start at its name on the left and follow the lines in the direction
of the arrows.  Every path through the diagram is a correct form of the
construct the diagram names.

\begin{itemize}
\item A word in a rectangle names another diagram, or one of the token
  classes defined in this chapter: \textit{identifier}, \textit{unsigned
  integer} and \textit{string literal}.
\item A reserved word in a rounded box, such as \textbf{begin}, is written
  exactly as shown (in capitals or small letters).
\item A symbol in a circle or a rounded box, such as \code{:=}, is written
  exactly as shown.
\item Where a line divides, choose any one branch.  A line that passes
  above an item may skip it; a line that returns below an item may repeat
  it.
\end{itemize}

The diagrams in this book are the grammar of the compiler.  Each one is
drawn from the rule that the Turbo Initech parser follows, and the
programs in this book, the compiler's own source and its test programs are
all derived from them.

\section{Character set and blanks}
\ix{character set}
\ix{blanks}
Turbo Initech programs are written in the printable ASCII characters.  The
space, the tab (character 9) and the characters 10 through 13 (line feed,
vertical tab, form feed and carriage return) are \term{blanks}: they
separate tokens and are otherwise ignored.  A blank may not appear inside a
token, except inside a character string.  Any other character that is not
part of a token or a comment is an error (``bad character,'' Appendix A).

Lines may be of any length; the whole source file may hold up to 262,144
characters (the limit \code{SourceMax} of the compiler, described in the
\doctitle{User's Guide}, Chapter 8).

\section{Special symbols}
\ix{special symbols}
\ix{operators}
Turbo Initech uses these single characters as special symbols:

\begin{BorTable}{Single-character symbols}{p{20mm}p{78mm}}
\Hd{Symbol} & \Hd{Use}\\\hline
\code{+ -} & addition and subtraction; unary minus\\
\code{*} & multiplication\\
\code{= < >} & relational operators; \code{=} also in declarations\\
\code{( )} & grouping, parameter lists and argument lists\\
\code{[ ]} & array bounds, string capacities and indexes\\
\code{.} & field selection; the end of the program\\
\code{,} & list separator\\
\code{:} & declarations of variables, parameters and results\\
\code{;} & separator of statements and declarations
\end{BorTable}

and these pairs of characters, which are single tokens and are written
without a space between the characters:

\begin{BorTable}{Two-character symbols}{p{20mm}p{78mm}}
\Hd{Symbol} & \Hd{Use}\\\hline
\code{:=} & assignment\\
\code{<> <= >=} & relational operators\\
\code{..} & the range of array bounds
\end{BorTable}

The apostrophe encloses character strings, and the braces \code{\{ \}} and
the pairs \code{(* *)} enclose comments.  Turbo Initech has no other
special symbols: there is no division operator \code{/}, and the
characters \code{\^} \code{@} \code{\$} and \code{\#} are not part of the
language.

\section{Reserved words}
\ix{reserved words}
The words in Table~\ref{tab:reserved} are \term{reserved words}.  They
have fixed meanings and cannot be used as identifiers.  Reserved words are
not case-sensitive: \code{BEGIN}, \code{Begin} and \code{begin} are the same
word.  In this book reserved words are printed in lower case.

\begin{BorTable}{Reserved words\label{tab:reserved}}{p{24mm}p{24mm}p{24mm}p{24mm}}
\ReservedItem{and} & \ReservedItem{downto} & \ReservedItem{length} & \ReservedItem{string}\\
\ReservedItem{array} & \ReservedItem{else} & \ReservedItem{mod} & \ReservedItem{then}\\
\ReservedItem{begin} & \ReservedItem{end} & \ReservedItem{not} & \ReservedItem{to}\\
\ReservedItem{boolean} & \ReservedItem{false} & \ReservedItem{of} & \ReservedItem{true}\\
\ReservedItem{case} & \ReservedItem{file} & \ReservedItem{or} & \ReservedItem{type}\\
\ReservedItem{char} & \ReservedItem{for} & \ReservedItem{ord} & \ReservedItem{until}\\
\ReservedItem{chr} & \ReservedItem{forward} & \ReservedItem{procedure} & \ReservedItem{var}\\
\ReservedItem{const} & \ReservedItem{function} & \ReservedItem{program} & \ReservedItem{while}\\
\ReservedItem{div} & \ReservedItem{if} & \ReservedItem{record} & \ReservedItem{write}\\
\ReservedItem{do} & \ReservedItem{integer} & \ReservedItem{repeat} & \ReservedItem{writeln}
\end{BorTable}

\MarginNote{Standard Pascal treats some of these words as predefined
identifiers.}Several of these words are reserved in Turbo Initech although
other Pascal compilers treat them as predefined identifiers that a program
may redeclare: the type names \rw{integer}, \rw{boolean} and \rw{char}, the
constants \rw{true} and \rw{false}, the directive \rw{forward}, and the
standard routines \rw{write}, \rw{writeln}, \rw{ord}, \rw{chr} and
\rw{length}.  The word \rw{case} is reserved, but Turbo Initech has no
\rw{case} statement; a program that uses the word is rejected (Chapter 6).

The names of the five file routines, \std{assign}{Assign},
\std{reset}{Reset}, \std{rewrite}{Rewrite}, \std{blockread}{BlockRead} and
\std{blockwrite}{BlockWrite}, are identifiers, not reserved words, but a
program may not declare anything with one of these names (Chapter 9).

\section{Identifiers}
\ix{identifiers}
\term{Identifiers} name the program, constants, types, variables, fields,
procedures, functions and parameters.

\Syntax{identifier}

A \textit{letter} is one of \code{A} through \code{Z} or \code{a} through
\code{z}; a \textit{digit} is one of \code{0} through \code{9}.  An
identifier may be up to 255 characters long, and all its characters are
significant.  Identifiers are not case-sensitive: the compiler converts every
identifier to small letters as it reads it, so \id{Total}, \id{TOTAL} and
\id{total} are the same identifier.  No identifier may be spelled like a
reserved word.

\section{Numbers}
\ix{numbers}
\ix{integer constants}
A number in a program is an \term{unsigned integer}, written in decimal:

\Syntax{unsigned-integer}

Its value must not exceed 2147483647, the largest value of the type
\rw{integer}.  There are no hexadecimal numbers and no real numbers.  A
minus sign before a number is the unary minus operator (Chapter 5), except
in constant declarations and array bounds, where the grammar allows it
directly.  The value $-2147483648$ cannot be written as a constant, since
2147483648 is not a valid number.

\section{Character strings}
\ix{strings, literal}
\ix{character strings}
A \term{string literal} is a sequence of characters enclosed in
apostrophes:

\Syntax{string-literal}

A \textit{string character} is any character except the apostrophe and the
line-end characters 10 and 13.  Two apostrophes in a row stand for one
apostrophe character in the string.  A string literal may hold up to 255
characters and may not continue past the end of a line.  The string with no
characters, \code{''}, is the empty string.

A string literal of exactly one character is a constant of type \rw{char};
any other string literal is a constant of type \rw{string} (Chapter 3).

\section{Comments}
\ix{comments}
A \term{comment} is ignored by the compiler, like a blank.

\Syntax{comment}

A comment begins with \code{\{} and ends with the next \code{\}}, or begins
with \code{(*} and ends with the next \code{*)}.  A comment may span any
number of lines.  Comments do not nest: inside a comment the characters
\code{\{} and \code{(*} have no meaning, and the first closing \code{\}} or
\code{*)} that matches the opening form ends the comment.  A comment that
begins with a dollar sign, \code{\{\$}, is an ordinary comment; Turbo
Initech has no compiler directives (Chapter 14).  A comment that is not
closed before the end of the source file is an error.

\section{Tokens in practice}
The program \file{TOKENS.PAS} uses each kind of token:

\Program{tokens}

It prints:

\ProgramOutput{tokens}

The assignment to \id{TOTAL} names the same variable as \id{Total} and
\id{total}.  Its value, $102 + 2147483600$, does not fit in an
\rw{integer}; Turbo Initech arithmetic keeps the low 32 bits of every
result, so the program prints $-2147483596$ (Chapter 5).

The compiler shows you the tokens it reads.  The first part of its output,
between the lines \code{TPS-LEX-BEGIN} and \code{TPS-LEX-END}, lists one token
per line.  For this program,

\Program{lexdemo}

the token list is:

\CompilerOutput{lexdemo}{LEX}

Reserved words are listed with \code{KW}, identifiers with \code{IDENT} and
their spelling in small letters, numbers with \code{INTEGER} and their value,
and strings with \code{STRING}, their length and the code of each character.
The comment produced no token.  The \doctitle{User's Guide}, Chapter 8,
describes the whole output of the compiler.
